\documentclass[11pt,a4paper]{article}

\usepackage[utf8]{inputenc}
\usepackage[T5]{fontenc}
\usepackage{amsmath,amssymb,amsthm,mathtools}
\usepackage{geometry}
\usepackage{enumitem}
\usepackage{hyperref}
\geometry{margin=2.4cm}

\newtheorem{theorem}{Định lý}
\newtheorem{lemma}{Bổ đề}
\newtheorem{proposition}{Mệnh đề}
\newtheorem{remark}{Nhận xét}
\newtheorem{assumption}{Giả thiết}
\newtheorem{corollary}{Hệ quả}
\newtheorem{definition}{Định nghĩa}

\newcommand{\erank}{\mathrm{erank}}
\newcommand{\tr}{\mathrm{tr}}
\newcommand{\diag}{\mathrm{diag}}
\newcommand{\RR}{\mathbb{R}}

\title{Ghi chú tiếng Việt về Theorem 2 trong \texttt{paper\_merged.tex}}
\author{}
\date{}

\begin{document}
\maketitle

\section*{Mục tiêu}
Tài liệu này làm ba việc:
\begin{enumerate}[nosep]
    \item viết lại và giải thích chi tiết Theorem 2 trong bản thảo;
    \item chỉ ra chính xác Theorem 2 có hay không phụ thuộc logic vào Theorem 1;
    \item đề xuất một phiên bản \emph{chứng minh được sạch hơn} của Theorem 2, với ít giả thiết hơn hoặc claim yếu hơn.
\end{enumerate}

\section*{Phát biểu gốc của Theorem 2}
Trong file \texttt{paper\_merged.tex}, Theorem 2 phát biểu:
\begin{quote}
Under Assumption~1, the preconditioned update
\[
W_{t+1}=W_t-\eta P_t^{-1}G_t
\]
with $\kappa(P)=\lambda_{\max}(P)/\lambda_{\min}(P)$ satisfies:
\begin{enumerate}[nosep]
    \item $\erank(P^{-1}G)\ge \erank(G)\min\!\bigl(1,\kappa(G)/\kappa(P)\bigr)$.
    \item When $\kappa(P)\le \kappa(G)$: $\erank(P^{-1}G)\ge \erank(G)$.
    \item $\erank(W_T)\ge \min\!\bigl(\erank(W_0),\, r(1-c/\kappa(P))\bigr)$ for some constant $c>0$.
\end{enumerate}
\end{quote}

\noindent
Điểm quan trọng là: \textbf{Assumption 1 trong paper chỉ nói về gradient concentration}, chứ chưa nói gì đủ mạnh về cách eigenbasis của $P$ ăn khớp với singular basis của $G$. Vì vậy, như phát biểu hiện tại, Theorem 2 còn thiếu một mắt xích kỹ thuật.

\section*{Kết luận ngắn trước khi vào chứng minh}
\begin{remark}
\leavevmode
\begin{enumerate}[nosep]
    \item \textbf{Theorem 2 không cần Theorem 1 về mặt logic.} Hai mệnh đề đầu của Theorem 2 là mệnh đề \emph{một bước} về toán tử $P^{-1}$ tác động lên gradient hiện tại $G$, nên có thể chứng minh trực tiếp mà không cần dùng Theorem 1.
    \item \textbf{Bản proof sketch hiện tại của paper có nhắc Theorem 1 ở bước cuối, nhưng đó là dùng như trực giác đối chiếu, không phải là phụ thuộc logic bắt buộc.}
    \item \textbf{Phần (3) của Theorem 2 hiện là phần yếu nhất}. Muốn chứng minh chặt, phải thêm giả thiết động lực học về việc các update đã precondition vẫn tiếp tục kích thích nhiều hướng và không bị triệt tiêu lẫn nhau.
\end{enumerate}
\end{remark}

\section*{Phiên bản có thể chứng minh được}
Để có một chứng minh kín, ta thêm giả thiết căn chỉnh phổ sau.

\begin{assumption}[Căn chỉnh phổ cục bộ]
\label{ass:alignment}
Xét một bước cập nhật cố định. Tồn tại các ma trận trực giao $U\in\RR^{m\times r}$, $V\in\RR^{n\times r}$ sao cho
\[
G = U \diag(s_1,\dots,s_r)V^\top,\qquad s_1\ge s_2\ge \cdots\ge s_r>0,
\]
và trên không gian con sinh bởi các hướng hoạt động của gradient, preconditioner tác động theo cùng basis trái:
\[
P\,U = U\diag(\lambda_1,\dots,\lambda_r),\qquad \lambda_1\ge \lambda_2\ge \cdots\ge \lambda_r>0.
\]
Nói cách khác, những hướng có năng lượng gradient lớn hơn cũng là những hướng có curvature lớn hơn, nên $P^{-1}$ sẽ giảm mạnh hơn ở các hướng đang trội.
\end{assumption}

\begin{remark}
Giả thiết này đúng chính xác trong mô hình quadratic đã được diagonal hóa, và là giả thiết chuẩn khi muốn nói ``preconditioning làm phẳng phổ cập nhật''. Nếu không có giả thiết này, chỉ từ Assumption 1 của paper chưa thể kết luận $\erank(P^{-1}G)\ge \erank(G)$, vì $P^{-1}$ hoàn toàn có thể khuếch đại nhầm một vài hướng vốn đã trội và làm phổ còn lệch hơn.
\end{remark}

\subsection*{Bước 1: viết singular values của update đã precondition}
Dưới Giả thiết~\ref{ass:alignment},
\[
P^{-1}G
=
U\diag(\lambda_1^{-1},\dots,\lambda_r^{-1})U^\top U\diag(s_1,\dots,s_r)V^\top
=
U\diag(t_1,\dots,t_r)V^\top,
\]
trong đó
\[
t_i = \frac{s_i}{\lambda_i},\qquad i=1,\dots,r.
\]
Như vậy singular values của update mới chính là singular values cũ được rescale từng hướng bởi $1/\lambda_i$.

\subsection*{Bước 2: condition number của update bị co lại}
\begin{lemma}[Co condition number]
\label{lem:kappa}
Dưới Giả thiết~\ref{ass:alignment},
\[
\kappa(P^{-1}G)=\frac{\kappa(G)}{\kappa(P)}.
\]
Do đó nếu $\kappa(P)>1$, độ lệch phổ của update giảm đi đúng một hệ số $\kappa(P)$.
\end{lemma}

\begin{proof}
Theo công thức trên,
\[
\kappa(P^{-1}G)=\frac{t_1}{t_r}
=
\frac{s_1/\lambda_1}{s_r/\lambda_r}
=
\frac{s_1}{s_r}\cdot \frac{\lambda_r}{\lambda_1}
=
\frac{\kappa(G)}{\kappa(P)}.
\]
Chứng minh xong.
\end{proof}

\subsection*{Bước 3: vì sao phổ được làm phẳng thì effective rank tăng}
Đặt
\[
p_i=\frac{s_i^2}{\sum_{j=1}^r s_j^2},
\qquad
\tilde p_i=\frac{t_i^2}{\sum_{j=1}^r t_j^2}
=
\frac{s_i^2/\lambda_i^2}{\sum_{j=1}^r s_j^2/\lambda_j^2}.
\]
Khi đó
\[
\erank(G)=\exp\!\Bigl(-\sum_{i=1}^r p_i\log p_i\Bigr),
\qquad
\erank(P^{-1}G)=\exp\!\Bigl(-\sum_{i=1}^r \tilde p_i\log \tilde p_i\Bigr).
\]

\noindent
Ý tưởng là rất đơn giản:
\begin{enumerate}[nosep]
    \item $s_1\ge \cdots\ge s_r$ nên năng lượng của $G$ đang tập trung mạnh hơn ở các chỉ số nhỏ.
    \item $\lambda_1\ge \cdots\ge \lambda_r$ nên $1/\lambda_1\le \cdots\le 1/\lambda_r$.
    \item Vì thế phép nhân bởi $P^{-1}$ giảm mạnh hơn ở các hướng đang trội và giảm yếu hơn ở các hướng đang yếu.
    \item Đây là đúng tinh thần của một phép \emph{Robin Hood transfer}: lấy bớt khối lượng ở đầu phổ và bù cho đuôi phổ.
    \item Entropy Shannon tăng dưới các phép làm đều như vậy; do đó effective rank tăng.
\end{enumerate}

\noindent
Để viết chặt điều này, ta dùng một giả thiết pairwise hơi mạnh hơn.

\begin{assumption}[Không đảo thứ tự sau preconditioning]
\label{ass:no-reorder}
Với mọi $i<j$,
\[
\frac{s_i^2}{s_j^2}\ge \frac{\lambda_i^2}{\lambda_j^2}.
\]
Tương đương với việc dãy $(t_i)$ vẫn không tăng:
\[
t_1\ge t_2\ge \cdots\ge t_r.
\]
\end{assumption}

\begin{lemma}[Majorization của phổ chuẩn hóa]
\label{lem:maj}
Dưới Giả thiết~\ref{ass:alignment} và~\ref{ass:no-reorder}, vector xác suất
$\tilde p=(\tilde p_1,\dots,\tilde p_r)$ bị majorize bởi $p=(p_1,\dots,p_r)$:
\[
\tilde p \prec p.
\]
Vì entropy Shannon là hàm Schur-concave, suy ra
\[
H(\tilde p)\ge H(p),
\qquad
\erank(P^{-1}G)\ge \erank(G).
\]
\end{lemma}

\begin{proof}
Ta có
\[
\frac{\tilde p_i}{\tilde p_j}
=
\frac{s_i^2/\lambda_i^2}{s_j^2/\lambda_j^2}
=
\frac{p_i}{p_j}\cdot \frac{\lambda_j^2}{\lambda_i^2}
\le
\frac{p_i}{p_j}
\qquad (i<j),
\]
vì $\lambda_i\ge \lambda_j$. Như vậy mọi tỉ lệ đầu/đuôi trong phổ chuẩn hóa đều bị co lại sau preconditioning. Dưới Giả thiết~\ref{ass:no-reorder}, thứ tự giảm vẫn được giữ nguyên, nên các tổng đầu
\[
\sum_{i=1}^k \tilde p_i
\]
không thể vượt quá
\[
\sum_{i=1}^k p_i,
\qquad k=1,\dots,r-1.
\]
Đây chính là định nghĩa của $\tilde p\prec p$. Từ đó
\[
H(\tilde p)\ge H(p).
\]
Lấy hàm mũ hai vế:
\[
\erank(P^{-1}G)=e^{H(\tilde p)}\ge e^{H(p)}=\erank(G).
\]
\end{proof}

\begin{corollary}[Phiên bản sạch của phần (2)]
\label{cor:part2}
Dưới Giả thiết~\ref{ass:alignment} và~\ref{ass:no-reorder},
\[
\erank(P^{-1}G)\ge \erank(G).
\]
\end{corollary}

\paragraph{Lưu ý quan trọng.}
Trong paper, phần (2) chỉ giả sử $\kappa(P)\le \kappa(G)$. Điều kiện này \emph{yếu hơn} Giả thiết~\ref{ass:no-reorder}; nó chỉ khống chế biên trên và biên dưới, chứ không khống chế mọi tỉ lệ cặp chỉ số. Vì vậy, nếu muốn giữ proof chặt, nên thay phần (2) bằng Corollary~\ref{cor:part2}, hoặc ít nhất phải nói rõ đó là một \emph{coarse sufficient condition under aligned spectra}.

\subsection*{Bước 4: một cận dưới định lượng chỉ dùng condition number}
Nếu muốn giữ thông điệp ``conditioning tốt hơn $\Rightarrow$ effective rank lớn hơn'' mà không cần so sánh trực tiếp với $\erank(G)$, ta có thể dùng một cận dưới chuẩn hơn.

\begin{lemma}[Cận dưới entropy theo condition number]
\label{lem:entropy-kappa}
Xét một ma trận hạng $r$ với singular values dương và condition number không quá $K\ge 1$. Khi đó
\[
\erank(X)\ge \phi_r(K),
\]
trong đó
\[
\phi_r(K)
=
\exp\!\Bigl(
-a\log a -(r-1)b\log b
\Bigr),
\]
với
\[
a=\frac{K^2}{K^2+r-1},
\qquad
b=\frac{1}{K^2+r-1}.
\]
\end{lemma}

\begin{proof}
Với condition number bị chặn bởi $K$, cấu hình entropy nhỏ nhất xảy ra khi một singular value đạt cực đại và $r-1$ singular values còn lại đạt cực tiểu, tức là phổ tỷ lệ với
\[
(K,1,\dots,1).
\]
Sau khi chuẩn hóa theo bình phương singular values, ta nhận được đúng các xác suất $a,b,\dots,b$. Entropy của cấu hình cực trị này là
\[
-a\log a -(r-1)b\log b.
\]
Vì mọi cấu hình khác ít lệch hơn đều có entropy lớn hơn, suy ra kết quả.
\end{proof}

\begin{corollary}[Hệ quả trực tiếp từ Bổ đề~\ref{lem:kappa}]
\label{cor:clean-bound}
Dưới Giả thiết~\ref{ass:alignment},
\[
\erank(P^{-1}G)\ge \phi_r\!\left(\frac{\kappa(G)}{\kappa(P)}\right).
\]
\end{corollary}

\noindent
Đây là phiên bản \emph{dễ chứng minh hơn} và về mặt kỹ thuật sạch hơn phần (1) trong paper. Nó không so sánh trực tiếp với $\erank(G)$, nhưng nói rất rõ rằng khi $\kappa(P)$ giảm thì cận dưới của effective rank của update tăng lên.

\section*{Về phần (3): rank của $W_T$ có thật sự suy ra không?}
Phần (3) trong Theorem 2 muốn đi từ mệnh đề một-bước về $P^{-1}G$ sang một cận dưới dài hạn cho $\erank(W_T)$. Bước này \textbf{không suy ra chỉ từ Assumption 1 cộng với hai bổ đề ở trên}. Lý do:
\begin{enumerate}[nosep]
    \item $W_T=W_0-\eta\sum_{t=0}^{T-1}P_t^{-1}G_t$ là tổng của nhiều ma trận;
    \item dù mỗi hạng tử có effective rank cao, các hạng tử vẫn có thể cùng dồn vào một vài hướng hoặc triệt tiêu nhau;
    \item muốn có cận dưới cho $\erank(W_T)$, phải thêm giả thiết động lực học về \emph{persistent excitation sau preconditioning}, và tốt nhất là cả một giả thiết kiểm soát sự đổi basis theo thời gian.
\end{enumerate}

\noindent
Một phiên bản đúng hơn là mệnh đề sau.

\begin{assumption}[Kích thích đều sau preconditioning]
\label{ass:persistent-precond}
Tồn tại một basis trực giao cố định $\{E_1,\dots,E_r\}$ của không gian hoạt động sao cho với mọi $t$,
\[
P_t^{-1}G_t = \sum_{i=1}^r \alpha_{t,i} E_i,
\]
và tồn tại các hằng số $0<\underline{a}\le \overline{a}<\infty$ sao cho
\[
\underline{a}
\le
\frac{\alpha_{t,i}^2}{\sum_{j=1}^r \alpha_{t,j}^2}
\le
\overline{a}
\qquad
\text{với mọi } i,t,
\]
đồng thời không có triệt tiêu dài hạn:
\[
\left\|\sum_{t=0}^{T-1}P_t^{-1}G_t\right\|_F^2
\ge
c_0\sum_{t=0}^{T-1}\|P_t^{-1}G_t\|_F^2
\]
cho một hằng số $c_0>0$ và mọi $T$ đủ lớn.
\end{assumption}

\begin{proposition}[Phiên bản dài hạn có thể chứng minh]
\label{prop:long-run}
Dưới Giả thiết~\ref{ass:persistent-precond}, tồn tại hằng số $c>0$ sao cho với $T$ đủ lớn,
\[
\erank(W_T)\ge c\,r.
\]
\end{proposition}

\begin{proof}[Ý tưởng chứng minh]
Viết
\[
W_T = W_0 - \eta\sum_{t=0}^{T-1}\sum_{i=1}^r \alpha_{t,i}E_i.
\]
Nhờ giả thiết kích thích đều, năng lượng tích lũy theo từng hướng $E_i$ đều cùng bậc với tổng năng lượng. Nhờ giả thiết không triệt tiêu dài hạn, chuẩn Frobenius của tổng không bị nhỏ đi quá mức so với tổng bình phương các bước cập nhật. Suy ra phân bố năng lượng của $W_T$ trên basis $\{E_i\}$ không thể tập trung vào hữu hạn rất nhỏ số hướng; do đó entropy của phân bố năng lượng được chặn dưới bởi một hằng số dương tỷ lệ với $\log r$, và vì thế $\erank(W_T)\ge cr$.
\end{proof}

\paragraph{Điểm mấu chốt.}
Mệnh đề trên \emph{không dùng Theorem 1}. Nó chỉ cần một giả thiết động lực học đủ trực tiếp cho bản thân chuỗi update đã precondition.

\section*{Trả lời trực tiếp ba câu hỏi của bạn}

\subsection*{1. Theorem 2 có cần Theorem 1 không?}
\textbf{Không cần về mặt logic.}

Chi tiết:
\begin{enumerate}[nosep]
    \item Phần (1) và phần (2) của Theorem 2 là mệnh đề một-bước về $P^{-1}G$ nên chứng minh trực tiếp.
    \item Phần (3) nếu muốn chứng minh chặt thì phải thêm giả thiết động lực học riêng; vẫn không cần Theorem 1.
    \item Theorem 1 chỉ hữu ích về \emph{mặt kể chuyện}: nó cho baseline first-order bị co rank, để rồi Theorem 2 nói second-order tránh được hiện tượng đó.
\end{enumerate}

\subsection*{2. Có thể lược bỏ Theorem 1 không?}
\textbf{Nếu mục tiêu chỉ là giữ lập luận cho Theorem 2 thì có thể lược bỏ.}

Nhưng về cấu trúc paper thì tôi không khuyên bỏ hoàn toàn, vì:
\begin{enumerate}[nosep]
    \item Proposition về first-order methods ở phần sau đang trỏ lại Theorem 1.
    \item Theorem 1 giải thích rất rõ cơ chế đối chứng: first-order dồn năng lượng vào low-rank gradient subspace.
    \item Nếu bỏ Theorem 1, bạn phải thay bằng một proposition hoặc remark đủ mạnh để giữ narrative.
\end{enumerate}

\subsection*{3. Có thể hạ claim để dùng ít assumption hơn không?}
\textbf{Có, và đây là phương án tôi khuyên dùng.}

Phiên bản gọn và sạch hơn:
\begin{enumerate}[nosep]
    \item Giữ một theorem một-bước: \emph{preconditioning co condition number của update}:
    \[
    \kappa(P^{-1}G)\le \kappa(G)/\kappa(P)
    \]
    dưới giả thiết căn chỉnh phổ cục bộ.
    \item Đưa mệnh đề ``effective rank không giảm'' thành một corollary dưới giả thiết mạnh hơn về không đảo thứ tự phổ.
    \item Tách phần dài hạn của $W_T$ thành một proposition riêng, với giả thiết \emph{persistent excitation after preconditioning}.
\end{enumerate}

\noindent
Làm như vậy sẽ có ba lợi ích:
\begin{enumerate}[nosep]
    \item chứng minh sạch hơn;
    \item không cần dùng Assumption 1 quá tải cho cả động lực học dài hạn lẫn mệnh đề phổ một-bước;
    \item giảm nguy cơ referee bắt lỗi ở bước ``each update has higher rank $\Rightarrow$ accumulated weight has high rank''.
\end{enumerate}

\section*{Khuyến nghị chỉnh sửa trong paper}
Nếu muốn sửa tối thiểu mà vẫn an toàn về mặt toán, tôi đề nghị:
\begin{enumerate}[nosep]
    \item Giữ Theorem 1 làm baseline.
    \item Sửa Theorem 2 thành theorem một-bước, chỉ nói về $P^{-1}G$.
    \item Đưa phần $\erank(W_T)\ge \min(\erank(W_0),r(1-c/\kappa(P)))$ xuống thành một proposition riêng với giả thiết bổ sung kiểu Giả thiết~\ref{ass:persistent-precond}.
    \item Ở appendix, nói rõ rằng cần spectral alignment hoặc local quadratic model để biện minh việc $P$ và $G$ có thể phân tích cùng basis.
\end{enumerate}

\end{document}
